\section{Aim and Objectives}\label{sec:aim-and-objectives}

The aim of this dissertation is to design, implement and evaluate a guarded
intent-recognition engine that constrains the raw output of a small, offline,
instruction-tuned language model with deterministic validation, so that the
combined system can act as a reliable head of an automated security pipeline in a
CPU-only, single-device deployment. To achieve this aim the following objectives
are pursued:

\begin{enumerate}
\item To design and implement the Intent Recognition Engine (IRE): a component
that receives free-text security requests and emits strictly
schema-conformant, authorised intent contracts, consuming a small offline
language model as a proposition generator whose output is validated
deterministically before it is handed onward.
\item To construct a purpose-built evaluation data set of security prompts,
covering well-formed requests and deliberately flawed prompts across the flaw
categories that matter in practice (adversarial phrasing, out-of-scope targets,
ambiguous or contradictory requests, fabricated parameters and identifiers, and
role violations), each annotated with gold labels.
\item To establish an evaluation protocol that quantifies, on the committed
experiment record, how much the guarded pipeline improves structured intent
recognition and flaw catching over raw model output, and that abates each guard
layer independently to estimate its individual contribution.
\item To compare two offline language-model backends under the identical
guardrail stack, in order to measure how much of the achieved safety is
backend-independent rather than an artefact of one model.
\end{enumerate}

Each objective is addressed by a corresponding chapter of this dissertation:
the design is developed in Chapter~\ref{ch:product-design}, the data set and protocol in
Chapter~\ref{ch:methodology}, the implementation in
Chapter~\ref{ch:implementation-testing}, and the evaluation and backend comparison in
Chapter~\ref{ch:results}.